% Options for packages loaded elsewhere
\PassOptionsToPackage{unicode}{hyperref}
\PassOptionsToPackage{hyphens}{url}
\PassOptionsToPackage{dvipsnames,svgnames,x11names}{xcolor}
%
\documentclass[
  letterpaper,
  DIV=11,
  numbers=noendperiod]{scrreprt}

\usepackage{amsmath,amssymb}
\usepackage{iftex}
\ifPDFTeX
  \usepackage[T1]{fontenc}
  \usepackage[utf8]{inputenc}
  \usepackage{textcomp} % provide euro and other symbols
\else % if luatex or xetex
  \usepackage{unicode-math}
  \defaultfontfeatures{Scale=MatchLowercase}
  \defaultfontfeatures[\rmfamily]{Ligatures=TeX,Scale=1}
\fi
\usepackage{lmodern}
\ifPDFTeX\else  
    % xetex/luatex font selection
\fi
% Use upquote if available, for straight quotes in verbatim environments
\IfFileExists{upquote.sty}{\usepackage{upquote}}{}
\IfFileExists{microtype.sty}{% use microtype if available
  \usepackage[]{microtype}
  \UseMicrotypeSet[protrusion]{basicmath} % disable protrusion for tt fonts
}{}
\makeatletter
\@ifundefined{KOMAClassName}{% if non-KOMA class
  \IfFileExists{parskip.sty}{%
    \usepackage{parskip}
  }{% else
    \setlength{\parindent}{0pt}
    \setlength{\parskip}{6pt plus 2pt minus 1pt}}
}{% if KOMA class
  \KOMAoptions{parskip=half}}
\makeatother
\usepackage{xcolor}
\setlength{\emergencystretch}{3em} % prevent overfull lines
\setcounter{secnumdepth}{5}
% Make \paragraph and \subparagraph free-standing
\ifx\paragraph\undefined\else
  \let\oldparagraph\paragraph
  \renewcommand{\paragraph}[1]{\oldparagraph{#1}\mbox{}}
\fi
\ifx\subparagraph\undefined\else
  \let\oldsubparagraph\subparagraph
  \renewcommand{\subparagraph}[1]{\oldsubparagraph{#1}\mbox{}}
\fi

\usepackage{color}
\usepackage{fancyvrb}
\newcommand{\VerbBar}{|}
\newcommand{\VERB}{\Verb[commandchars=\\\{\}]}
\DefineVerbatimEnvironment{Highlighting}{Verbatim}{commandchars=\\\{\}}
% Add ',fontsize=\small' for more characters per line
\usepackage{framed}
\definecolor{shadecolor}{RGB}{241,243,245}
\newenvironment{Shaded}{\begin{snugshade}}{\end{snugshade}}
\newcommand{\AlertTok}[1]{\textcolor[rgb]{0.68,0.00,0.00}{#1}}
\newcommand{\AnnotationTok}[1]{\textcolor[rgb]{0.37,0.37,0.37}{#1}}
\newcommand{\AttributeTok}[1]{\textcolor[rgb]{0.40,0.45,0.13}{#1}}
\newcommand{\BaseNTok}[1]{\textcolor[rgb]{0.68,0.00,0.00}{#1}}
\newcommand{\BuiltInTok}[1]{\textcolor[rgb]{0.00,0.23,0.31}{#1}}
\newcommand{\CharTok}[1]{\textcolor[rgb]{0.13,0.47,0.30}{#1}}
\newcommand{\CommentTok}[1]{\textcolor[rgb]{0.37,0.37,0.37}{#1}}
\newcommand{\CommentVarTok}[1]{\textcolor[rgb]{0.37,0.37,0.37}{\textit{#1}}}
\newcommand{\ConstantTok}[1]{\textcolor[rgb]{0.56,0.35,0.01}{#1}}
\newcommand{\ControlFlowTok}[1]{\textcolor[rgb]{0.00,0.23,0.31}{#1}}
\newcommand{\DataTypeTok}[1]{\textcolor[rgb]{0.68,0.00,0.00}{#1}}
\newcommand{\DecValTok}[1]{\textcolor[rgb]{0.68,0.00,0.00}{#1}}
\newcommand{\DocumentationTok}[1]{\textcolor[rgb]{0.37,0.37,0.37}{\textit{#1}}}
\newcommand{\ErrorTok}[1]{\textcolor[rgb]{0.68,0.00,0.00}{#1}}
\newcommand{\ExtensionTok}[1]{\textcolor[rgb]{0.00,0.23,0.31}{#1}}
\newcommand{\FloatTok}[1]{\textcolor[rgb]{0.68,0.00,0.00}{#1}}
\newcommand{\FunctionTok}[1]{\textcolor[rgb]{0.28,0.35,0.67}{#1}}
\newcommand{\ImportTok}[1]{\textcolor[rgb]{0.00,0.46,0.62}{#1}}
\newcommand{\InformationTok}[1]{\textcolor[rgb]{0.37,0.37,0.37}{#1}}
\newcommand{\KeywordTok}[1]{\textcolor[rgb]{0.00,0.23,0.31}{#1}}
\newcommand{\NormalTok}[1]{\textcolor[rgb]{0.00,0.23,0.31}{#1}}
\newcommand{\OperatorTok}[1]{\textcolor[rgb]{0.37,0.37,0.37}{#1}}
\newcommand{\OtherTok}[1]{\textcolor[rgb]{0.00,0.23,0.31}{#1}}
\newcommand{\PreprocessorTok}[1]{\textcolor[rgb]{0.68,0.00,0.00}{#1}}
\newcommand{\RegionMarkerTok}[1]{\textcolor[rgb]{0.00,0.23,0.31}{#1}}
\newcommand{\SpecialCharTok}[1]{\textcolor[rgb]{0.37,0.37,0.37}{#1}}
\newcommand{\SpecialStringTok}[1]{\textcolor[rgb]{0.13,0.47,0.30}{#1}}
\newcommand{\StringTok}[1]{\textcolor[rgb]{0.13,0.47,0.30}{#1}}
\newcommand{\VariableTok}[1]{\textcolor[rgb]{0.07,0.07,0.07}{#1}}
\newcommand{\VerbatimStringTok}[1]{\textcolor[rgb]{0.13,0.47,0.30}{#1}}
\newcommand{\WarningTok}[1]{\textcolor[rgb]{0.37,0.37,0.37}{\textit{#1}}}

\providecommand{\tightlist}{%
  \setlength{\itemsep}{0pt}\setlength{\parskip}{0pt}}\usepackage{longtable,booktabs,array}
\usepackage{calc} % for calculating minipage widths
% Correct order of tables after \paragraph or \subparagraph
\usepackage{etoolbox}
\makeatletter
\patchcmd\longtable{\par}{\if@noskipsec\mbox{}\fi\par}{}{}
\makeatother
% Allow footnotes in longtable head/foot
\IfFileExists{footnotehyper.sty}{\usepackage{footnotehyper}}{\usepackage{footnote}}
\makesavenoteenv{longtable}
\usepackage{graphicx}
\makeatletter
\def\maxwidth{\ifdim\Gin@nat@width>\linewidth\linewidth\else\Gin@nat@width\fi}
\def\maxheight{\ifdim\Gin@nat@height>\textheight\textheight\else\Gin@nat@height\fi}
\makeatother
% Scale images if necessary, so that they will not overflow the page
% margins by default, and it is still possible to overwrite the defaults
% using explicit options in \includegraphics[width, height, ...]{}
\setkeys{Gin}{width=\maxwidth,height=\maxheight,keepaspectratio}
% Set default figure placement to htbp
\makeatletter
\def\fps@figure{htbp}
\makeatother
% definitions for citeproc citations
\NewDocumentCommand\citeproctext{}{}
\NewDocumentCommand\citeproc{mm}{%
  \begingroup\def\citeproctext{#2}\cite{#1}\endgroup}
\makeatletter
 % allow citations to break across lines
 \let\@cite@ofmt\@firstofone
 % avoid brackets around text for \cite:
 \def\@biblabel#1{}
 \def\@cite#1#2{{#1\if@tempswa , #2\fi}}
\makeatother
\newlength{\cslhangindent}
\setlength{\cslhangindent}{1.5em}
\newlength{\csllabelwidth}
\setlength{\csllabelwidth}{3em}
\newenvironment{CSLReferences}[2] % #1 hanging-indent, #2 entry-spacing
 {\begin{list}{}{%
  \setlength{\itemindent}{0pt}
  \setlength{\leftmargin}{0pt}
  \setlength{\parsep}{0pt}
  % turn on hanging indent if param 1 is 1
  \ifodd #1
   \setlength{\leftmargin}{\cslhangindent}
   \setlength{\itemindent}{-1\cslhangindent}
  \fi
  % set entry spacing
  \setlength{\itemsep}{#2\baselineskip}}}
 {\end{list}}
\usepackage{calc}
\newcommand{\CSLBlock}[1]{\hfill\break\parbox[t]{\linewidth}{\strut\ignorespaces#1\strut}}
\newcommand{\CSLLeftMargin}[1]{\parbox[t]{\csllabelwidth}{\strut#1\strut}}
\newcommand{\CSLRightInline}[1]{\parbox[t]{\linewidth - \csllabelwidth}{\strut#1\strut}}
\newcommand{\CSLIndent}[1]{\hspace{\cslhangindent}#1}

\KOMAoption{captions}{tableheading}
\makeatletter
\@ifpackageloaded{tcolorbox}{}{\usepackage[skins,breakable]{tcolorbox}}
\@ifpackageloaded{fontawesome5}{}{\usepackage{fontawesome5}}
\definecolor{quarto-callout-color}{HTML}{909090}
\definecolor{quarto-callout-note-color}{HTML}{0758E5}
\definecolor{quarto-callout-important-color}{HTML}{CC1914}
\definecolor{quarto-callout-warning-color}{HTML}{EB9113}
\definecolor{quarto-callout-tip-color}{HTML}{00A047}
\definecolor{quarto-callout-caution-color}{HTML}{FC5300}
\definecolor{quarto-callout-color-frame}{HTML}{acacac}
\definecolor{quarto-callout-note-color-frame}{HTML}{4582ec}
\definecolor{quarto-callout-important-color-frame}{HTML}{d9534f}
\definecolor{quarto-callout-warning-color-frame}{HTML}{f0ad4e}
\definecolor{quarto-callout-tip-color-frame}{HTML}{02b875}
\definecolor{quarto-callout-caution-color-frame}{HTML}{fd7e14}
\makeatother
\makeatletter
\@ifpackageloaded{bookmark}{}{\usepackage{bookmark}}
\makeatother
\makeatletter
\@ifpackageloaded{caption}{}{\usepackage{caption}}
\AtBeginDocument{%
\ifdefined\contentsname
  \renewcommand*\contentsname{Table of contents}
\else
  \newcommand\contentsname{Table of contents}
\fi
\ifdefined\listfigurename
  \renewcommand*\listfigurename{List of Figures}
\else
  \newcommand\listfigurename{List of Figures}
\fi
\ifdefined\listtablename
  \renewcommand*\listtablename{List of Tables}
\else
  \newcommand\listtablename{List of Tables}
\fi
\ifdefined\figurename
  \renewcommand*\figurename{Figure}
\else
  \newcommand\figurename{Figure}
\fi
\ifdefined\tablename
  \renewcommand*\tablename{Table}
\else
  \newcommand\tablename{Table}
\fi
}
\@ifpackageloaded{float}{}{\usepackage{float}}
\floatstyle{ruled}
\@ifundefined{c@chapter}{\newfloat{codelisting}{h}{lop}}{\newfloat{codelisting}{h}{lop}[chapter]}
\floatname{codelisting}{Listing}
\newcommand*\listoflistings{\listof{codelisting}{List of Listings}}
\makeatother
\makeatletter
\makeatother
\makeatletter
\@ifpackageloaded{caption}{}{\usepackage{caption}}
\@ifpackageloaded{subcaption}{}{\usepackage{subcaption}}
\makeatother
\ifLuaTeX
  \usepackage{selnolig}  % disable illegal ligatures
\fi
\usepackage{bookmark}

\IfFileExists{xurl.sty}{\usepackage{xurl}}{} % add URL line breaks if available
\urlstyle{same} % disable monospaced font for URLs
\hypersetup{
  pdftitle={Notes de cours - MAT8186},
  pdfauthor={Cedric Beaulac},
  colorlinks=true,
  linkcolor={blue},
  filecolor={Maroon},
  citecolor={Blue},
  urlcolor={Blue},
  pdfcreator={LaTeX via pandoc}}

\title{Notes de cours - MAT8186}
\author{Cedric Beaulac}
\date{2026-08-01}

\begin{document}
\maketitle

\renewcommand*\contentsname{Table of contents}
{
\hypersetup{linkcolor=}
\setcounter{tocdepth}{2}
\tableofcontents
}
\bookmarksetup{startatroot}

\chapter*{Index}\label{index}
\addcontentsline{toc}{chapter}{Index}

\markboth{Index}{Index}

Il s'agit d'un livre Quarto contenant le matériel du cours MAT8186.

Ce matériel fut conçu pour un cours de 15 heures organisé en 5 séances
de 3 heures. L'objectif du cours est d'amener les étudiants à développer
leur propres librairies R et s'inspire fortement du livre
\href{https://r-pkgs.org}{R Packages}.

\bookmarksetup{startatroot}

\chapter{Un survol rapide}\label{un-survol-rapide}

Dans ce cours, nous apprendrons comment mettre au point une librairie R.

\section*{Pourquoi ?}\label{pourquoi}
\addcontentsline{toc}{section}{Pourquoi ?}

\markright{Pourquoi ?}

La librairie (de l'anglais \emph{package}) est le format le mieux
structuré pour le partage de code. Aux cycles supérieurs, les étudiants
sont appelés à entrer en contact avec la recherche scientifique. Une
étape importante de la recherche est la dissémination des résultats.
Dans le cadre de la recherche en science des données, il est commun de
développer de nouvelles méthodologies d'analyse de données et
d'implémenter celles-ci. La mise au point d'une librairie est donc une
manière de diffuser ces implémentations et, conséquemment, de rendre de
nouvelles méthodes d'analyse de données facilement accessibles.

L'objectif est donc de remplacer le partage d'un simple document
\texttt{.R} par un ensemble de fonctions réfléchies, complètes,
documentées et facilement accessibles à de nouveaux utilisateurs.

L'ajout constant de nouvelles librairies R sur le \emph{Comprehensive R
Archive Network} (CRAN) est l'une des raisons du succès de R au fil des
années et permet à R de rester un langage utilisé et reconnu à l'ère de
Python, Julia et des autres langages couramment utilisés en science des
données. CRAN contient aujourd'hui près de 20 000 librairies! Vous avez
sûrement déjà installé des librairies à partir de CRAN avec la commande
\texttt{install.packages("x")}.

\section*{Pourquoi moi ?}\label{pourquoi-moi}
\addcontentsline{toc}{section}{Pourquoi moi ?}

\markright{Pourquoi moi ?}

Je ne suis pas un expert de la mise au point de librairies, mais :

\begin{itemize}
\tightlist
\item
  J'ai déjà mis au point une librairie!
\item
  Je veux mettre au point d'autres librairies!
\item
  J'aime R, l'enseignement et l'enseignement de R.
\item
  Personne d'autre ne pouvait l'enseigner.
\item
  Je compte sur une excellente référence pour faire une partie du
  travail à ma place.
\end{itemize}

\section*{Quoi ?}\label{quoi}
\addcontentsline{toc}{section}{Quoi ?}

\markright{Quoi ?}

Dans ce cours, nous allons bien entendu mettre au point une librairie R
: c'est l'objectif principal. Nous verrons comment passer d'un ensemble
de fonctions R à une librairie structurée, documentée, testée et pouvant
être facilement partagée avec d'autres utilisateurs.

Nous verrons également plusieurs outils et concepts qui accompagnent le
développement moderne en R :

\begin{itemize}
\tightlist
\item
  \textbf{Git et GitHub}, pour conserver un historique des différentes
  versions de notre code, partager celui-ci et collaborer avec d'autres
  utilisateurs;
\item
  les bonnes pratiques pour organiser des \textbf{projets, fichiers et
  fonctions R};
\item
  la \textbf{documentation et les tests}, afin qu'une librairie puisse
  être comprise, utilisée et modifiée de manière fiable;
\item
  quelques notions de \textbf{programmation orientée objet}, afin de
  comprendre comment R permet de définir des comportements associés à
  différents types d'objets;
\item
  \textbf{Rcpp et reticulate}, qui permettent respectivement d'intégrer
  du code C++ et Python à un projet R, notamment lorsque certaines
  fonctionnalités sont plus faciles ou plus efficaces à implémenter dans
  un autre langage.
\end{itemize}

Ces différents outils ne constituent pas des objectifs indépendants :
nous les utiliserons progressivement comme des composantes du
développement de notre librairie.

\section*{Comment ?}\label{comment}
\addcontentsline{toc}{section}{Comment ?}

\markright{Comment ?}

Nous allons essentiellement suivre le livre \href{https://r-pkgs.org}{R
Packages} de Hadley Wickham et Jennifer Bryan. Wickham est notamment
l'auteur ou le coauteur de plusieurs librairies importantes de
l'écosystème R, dont \emph{ggplot2}, et l'un des principaux
contributeurs au \emph{tidyverse} et aux outils modernes de
développement de librairies R.

Pour la mise au point de librairies, nous utiliserons\ldots{} une
librairie.

\section{Le concept de A à Z}\label{le-concept-de-a-uxe0-z}

Dans cette première partie du cours, nous allons parcourir \textbf{très
rapidement l'ensemble du processus de mise au point d'une librairie R}.

L'objectif n'est pas de comprendre immédiatement toutes les commandes,
tous les fichiers et toutes les conventions que nous allons rencontrer.

Nous allons plutôt commencer par voir \textbf{le chemin de A à Z}.

À la fin, nous aurons :

\begin{enumerate}
\def\labelenumi{\arabic{enumi}.}
\tightlist
\item
  créé une librairie R;
\item
  créé un dépôt Git;
\item
  ajouté une première fonction;
\item
  chargé notre librairie pendant son développement;
\item
  vérifié sa structure;
\item
  ajouté une licence;
\item
  documenté notre fonction;
\item
  ajouté un test;
\item
  déclaré une dépendance;
\item
  installé notre librairie;
\item
  vu comment elle pourra éventuellement être partagée.
\end{enumerate}

\begin{tcolorbox}[enhanced jigsaw, arc=.35mm, title=\textcolor{quarto-callout-note-color}{\faInfo}\hspace{0.5em}{Une première visite}, colback=white, left=2mm, rightrule=.15mm, toptitle=1mm, opacityback=0, bottomtitle=1mm, breakable, colframe=quarto-callout-note-color-frame, coltitle=black, toprule=.15mm, titlerule=0mm, opacitybacktitle=0.6, bottomrule=.15mm, colbacktitle=quarto-callout-note-color!10!white, leftrule=.75mm]

Il n'est \textbf{pas attendu que vous reteniez toutes les commandes
présentées dans cette section}.

Nous allons revenir beaucoup plus en détail sur chacune des étapes au
cours des prochaines séances.

\end{tcolorbox}

\section*{Nos principaux outils}\label{nos-principaux-outils}
\addcontentsline{toc}{section}{Nos principaux outils}

\markright{Nos principaux outils}

R possède déjà les outils nécessaires pour construire, installer et
vérifier des librairies.

Cependant, plusieurs librairies facilitent énormément le travail de
développement.

Nous utiliserons principalement \textbf{\texttt{devtools}} et
\textbf{\texttt{usethis}}.

\begin{Shaded}
\begin{Highlighting}[]
\FunctionTok{install.packages}\NormalTok{(}\StringTok{"devtools"}\NormalTok{)}
\end{Highlighting}
\end{Shaded}

Une fois installée :

\begin{Shaded}
\begin{Highlighting}[]
\FunctionTok{library}\NormalTok{(devtools)}
\FunctionTok{packageVersion}\NormalTok{(}\StringTok{"devtools"}\NormalTok{)}
\end{Highlighting}
\end{Shaded}

La librairie \texttt{devtools} rassemble plusieurs outils associés au
développement de librairies R. Le manuel de référence et la libraire
sont conçu pour être utilisé main dans la main. Cette libraire permet à
quiconque de \emph{facilement} mettre des librairies au point.

Une partie des fonctions que nous utiliserons proviennent en réalité de
la librairie \texttt{usethis}.

Par exemple :

\begin{Shaded}
\begin{Highlighting}[]
\NormalTok{usethis}\SpecialCharTok{::}\FunctionTok{create\_package}\NormalTok{()}
\end{Highlighting}
\end{Shaded}

Dans les notes de cours, nous utiliserons régulièrement la notation

\begin{Shaded}
\begin{Highlighting}[]
\NormalTok{librairie::fonction()}
\end{Highlighting}
\end{Shaded}

afin d'identifier clairement la provenance d'une fonction.

\section{Créons une librairie}\label{cruxe9ons-une-librairie}

Nous allons développer une petite librairie appelée :

\begin{Shaded}
\begin{Highlighting}[]
\NormalTok{mat8186demo}
\end{Highlighting}
\end{Shaded}

Elle ne fera rien de particulièrement révolutionnaire.

Son objectif sera simplement de nous permettre de parcourir une première
fois toutes les étapes du développement d'une librairie R.

Comme d'habitude, commencez par crée un répertoire et l'identifier comme
répertoir de travail:

\begin{Shaded}
\begin{Highlighting}[]
\FunctionTok{setwd}\NormalTok{(}\StringTok{"C:/VotreRepertoire"}\NormalTok{)}
\end{Highlighting}
\end{Shaded}

\section*{Création}\label{cruxe9ation}
\addcontentsline{toc}{section}{Création}

\markright{Création}

Puis, nous pouvons créer notre librairie :

\begin{Shaded}
\begin{Highlighting}[]
\NormalTok{usethis}\SpecialCharTok{::}\FunctionTok{create\_package}\NormalTok{(}\StringTok{"\textasciitilde{}/mat8186demo"}\NormalTok{)}
\end{Highlighting}
\end{Shaded}

\texttt{create\_package()} crée automatiquement plusieurs fichiers et
répertoires.

Nous obtenons quelque chose qui ressemble à :

\begin{Shaded}
\begin{Highlighting}[]
\NormalTok{mat8186demo/}
\NormalTok{├── .Rbuildignore}
\NormalTok{├── .gitignore}
\NormalTok{├── DESCRIPTION}
\NormalTok{├── NAMESPACE}
\NormalTok{├── R/}
\NormalTok{└── mat8186demo.Rproj}
\end{Highlighting}
\end{Shaded}

Nous venons déjà de créer plusieurs choses à la fois.

Notre dossier est maintenant :

\begin{itemize}
\tightlist
\item
  un \textbf{répertoire} sur notre ordinateur;
\item
  un \textbf{projet RStudio};
\item
  la structure minimale d'une \textbf{librairie R}.
\end{itemize}

\begin{tcolorbox}[enhanced jigsaw, arc=.35mm, title=\textcolor{quarto-callout-tip-color}{\faLightbulb}\hspace{0.5em}{Ne cherchons pas encore à tout comprendre}, colback=white, left=2mm, rightrule=.15mm, toptitle=1mm, opacityback=0, bottomtitle=1mm, breakable, colframe=quarto-callout-tip-color-frame, coltitle=black, toprule=.15mm, titlerule=0mm, opacitybacktitle=0.6, bottomrule=.15mm, colbacktitle=quarto-callout-tip-color!10!white, leftrule=.75mm]

Nous reviendrons plus tard sur \texttt{DESCRIPTION}, \texttt{NAMESPACE},
\texttt{.Rbuildignore} et les autres fichiers.

Pour l'instant, retenons surtout que le répertoire \texttt{R/}
contiendra le code de notre librairie.

\end{tcolorbox}

\section{Ajoutons Git}\label{ajoutons-git}

Avant même d'écrire notre première fonction, nous pouvons commencer à
conserver un historique de notre travail.

\begin{Shaded}
\begin{Highlighting}[]
\NormalTok{usethis}\SpecialCharTok{::}\FunctionTok{use\_git}\NormalTok{()}
\end{Highlighting}
\end{Shaded}

Notre même dossier est maintenant :

\begin{Shaded}
\begin{Highlighting}[]
\NormalTok{un répertoire}
\NormalTok{      +}
\NormalTok{un projet RStudio}
\NormalTok{      +}
\NormalTok{une librairie R}
\NormalTok{      +}
\NormalTok{un dépôt Git}
\end{Highlighting}
\end{Shaded}

Git peut maintenant conserver les différentes versions de notre travail.

Nous pourrions immédiatement enregistrer cette première étape avec un
\emph{commit} :

\begin{Shaded}
\begin{Highlighting}[]
\NormalTok{Initial package structure}
\end{Highlighting}
\end{Shaded}

Nous reviendrons beaucoup plus en détail sur \textbf{Git et GitHub} plus
tard dans cette séance.

Pour l'instant, Git conserve toute les version du répertoire nous
permettant de conserver une trace de toute les modifications faites.

\section{Une première fonction}\label{une-premiuxe8re-fonction}

Une librairie R contient principalement des \textbf{fonctions}. C'est
par un ensemble de fonction que l'on peut encoder des méthodes plus
complexes, l'atome d'une libraire est donc la fonction.

Créons une petite fonction permettant de centrer-réduire un vecteur
numérique.

Nous aimerions pouvoir faire :

\begin{Shaded}
\begin{Highlighting}[]
\FunctionTok{CR}\NormalTok{(}\FunctionTok{c}\NormalTok{(}\DecValTok{10}\NormalTok{, }\DecValTok{20}\NormalTok{, }\DecValTok{30}\NormalTok{))}
\end{Highlighting}
\end{Shaded}

et obtenir :

\begin{Shaded}
\begin{Highlighting}[]
\NormalTok{[1] {-}1 0 1}
\end{Highlighting}
\end{Shaded}

Une première version de notre fonction pourrait être :

\begin{Shaded}
\begin{Highlighting}[]
\NormalTok{CR }\OtherTok{\textless{}{-}} \ControlFlowTok{function}\NormalTok{(x) \{}
\NormalTok{  (x }\SpecialCharTok{{-}} \FunctionTok{mean}\NormalTok{(x)) }\SpecialCharTok{/} \FunctionTok{sd}\NormalTok{(x)}
\NormalTok{\}}
\end{Highlighting}
\end{Shaded}

Dans une librairie, notre code R est placé dans le répertoire :

\begin{Shaded}
\begin{Highlighting}[]
\NormalTok{R/}
\end{Highlighting}
\end{Shaded}

Nous pouvons créer un fichier pour notre fonction avec :

\begin{Shaded}
\begin{Highlighting}[]
\NormalTok{usethis}\SpecialCharTok{::}\FunctionTok{use\_r}\NormalTok{(}\StringTok{"CR"}\NormalTok{)}
\end{Highlighting}
\end{Shaded}

Cette commande crée :

\begin{Shaded}
\begin{Highlighting}[]
\NormalTok{R/CR.R}
\end{Highlighting}
\end{Shaded}

Nous pouvons alors y placer notre fonction :

\begin{Shaded}
\begin{Highlighting}[]
\NormalTok{CR }\OtherTok{\textless{}{-}} \ControlFlowTok{function}\NormalTok{(x) \{}
\NormalTok{  (x }\SpecialCharTok{{-}} \FunctionTok{mean}\NormalTok{(x)) }\SpecialCharTok{/} \FunctionTok{sd}\NormalTok{(x)}
\NormalTok{\}}
\end{Highlighting}
\end{Shaded}

\section*{Du script à la librairie}\label{du-script-uxe0-la-librairie}
\addcontentsline{toc}{section}{Du script à la librairie}

\markright{Du script à la librairie}

Il y a ici un changement important dans notre manière de travailler.

Dans un script d'analyse, nous pouvons retrouver quelque chose comme :

\begin{Shaded}
\begin{Highlighting}[]
\NormalTok{donnees }\OtherTok{\textless{}{-}} \FunctionTok{read.csv}\NormalTok{(}\StringTok{"donnees.csv"}\NormalTok{)}

\NormalTok{resultat }\OtherTok{\textless{}{-}} \FunctionTok{lm}\NormalTok{(y }\SpecialCharTok{\textasciitilde{}}\NormalTok{ x, }\AttributeTok{data =}\NormalTok{ donnees)}

\FunctionTok{summary}\NormalTok{(resultat)}

\FunctionTok{plot}\NormalTok{(resultat)}
\end{Highlighting}
\end{Shaded}

Dans le répertoire \texttt{R/} d'une librairie, nous retrouverons plutôt
des \textbf{définitions de fonctions} :

\begin{Shaded}
\begin{Highlighting}[]
\NormalTok{ma\_fonction }\OtherTok{\textless{}{-}} \ControlFlowTok{function}\NormalTok{(x) \{}
\NormalTok{  ...}
\NormalTok{\}}
\end{Highlighting}
\end{Shaded}

Nous passons progressivement d'une logique de \textbf{script d'analyse}
à une logique de \textbf{développement logiciel}.

\section{Utiliser notre librairie pendant son
développement}\label{utiliser-notre-librairie-pendant-son-duxe9veloppement}

Notre fonction est maintenant enregistrée dans :

\begin{Shaded}
\begin{Highlighting}[]
\NormalTok{R/CR.R}
\end{Highlighting}
\end{Shaded}

Comment pouvons-nous l'utiliser?

Nous pouvons faire :

\begin{Shaded}
\begin{Highlighting}[]
\FunctionTok{source}\NormalTok{(}\StringTok{"R/CR.R"}\NormalTok{)}
\end{Highlighting}
\end{Shaded}

et c'est en fait ce que nous faisons au premier cycle quand nous codons
nos propres fonctions. Ce n'est pas la façon privilégiée de travailler
lorsqu'on développe une librairie; nous allons \emph{installer} la
librairie.

Nous allons plutôt utiliser :

\begin{Shaded}
\begin{Highlighting}[]
\NormalTok{devtools}\SpecialCharTok{::}\FunctionTok{load\_all}\NormalTok{()}
\end{Highlighting}
\end{Shaded}

\texttt{load\_all()} charge temporairement notre librairie en cours de
développement.

Nous pouvons alors utiliser :

\begin{Shaded}
\begin{Highlighting}[]
\FunctionTok{CR}\NormalTok{(}\FunctionTok{c}\NormalTok{(}\DecValTok{10}\NormalTok{, }\DecValTok{20}\NormalTok{, }\DecValTok{30}\NormalTok{))}
\end{Highlighting}
\end{Shaded}

comme si notre librairie était déjà installée.

Nous obtenons notre premier \textbf{cycle de développement} :

\begin{Shaded}
\begin{Highlighting}[]
\NormalTok{modifier le code}
\NormalTok{       ↓}
\NormalTok{   load\_all()}
\NormalTok{       ↓}
\NormalTok{tester la fonction}
\NormalTok{       ↓}
\NormalTok{modifier le code}
\NormalTok{       ↓}
\NormalTok{   load\_all()}
\NormalTok{       ↓}
\NormalTok{      ...}
\end{Highlighting}
\end{Shaded}

Ce cycle reviendra constamment pendant le cours.

\section{Vérifier notre librairie}\label{vuxe9rifier-notre-librairie}

Une librairie R doit respecter plusieurs conventions, normes afin
d'opérer. Ces normes et standards sont la raison principale pourquoi
aurjourd'hui il est beaucoup plus facile d'utiliser des librairies R que
des paquet Python.

Heureusement, nous n'avons pas besoin de tous vérifier manuellement
grâce à \texttt{devtools}. Nous pouvons demander à R de vérifier notre
librairie avec :

\begin{Shaded}
\begin{Highlighting}[]
\NormalTok{devtools}\SpecialCharTok{::}\FunctionTok{check}\NormalTok{()}
\end{Highlighting}
\end{Shaded}

Cette commande réalise une série de vérifications sur notre librairie.

Le résultat peut contenir trois types de problèmes :

\begin{Shaded}
\begin{Highlighting}[]
\NormalTok{ERROR}
\NormalTok{WARNING}
\NormalTok{NOTE}
\end{Highlighting}
\end{Shaded}

À terme, nous aimerions obtenir :

\begin{Shaded}
\begin{Highlighting}[]
\NormalTok{0 errors ✔}
\NormalTok{0 warnings ✔}
\NormalTok{0 notes ✔}
\end{Highlighting}
\end{Shaded}

\begin{tcolorbox}[enhanced jigsaw, arc=.35mm, title=\textcolor{quarto-callout-important-color}{\faExclamation}\hspace{0.5em}{Vérifier régulièrement}, colback=white, left=2mm, rightrule=.15mm, toptitle=1mm, opacityback=0, bottomtitle=1mm, breakable, colframe=quarto-callout-important-color-frame, coltitle=black, toprule=.15mm, titlerule=0mm, opacitybacktitle=0.6, bottomrule=.15mm, colbacktitle=quarto-callout-important-color!10!white, leftrule=.75mm]

Il est beaucoup plus facile de corriger un problème lorsqu'il vient
d'apparaître que lorsqu'on a modifié vingt fichiers depuis notre
dernière vérification.

Nous utiliserons donc \texttt{check()} \textbf{régulièrement pendant le
développement}.

\end{tcolorbox}

\section{Le fichier DESCRIPTION}\label{le-fichier-description}

Notre librairie possède déjà un fichier nommé :

\begin{Shaded}
\begin{Highlighting}[]
\NormalTok{DESCRIPTION}
\end{Highlighting}
\end{Shaded}

Ce fichier contient les principales \textbf{métadonnées} de notre
librairie.

Il ressemble initialement à ceci :

\begin{Shaded}
\begin{Highlighting}[]
\NormalTok{Package: mat8186demo}
\NormalTok{Title: What the Package Does}
\NormalTok{Version: 0.0.0.9000}
\NormalTok{Authors@R: ...}
\NormalTok{Description: What the package does.}
\NormalTok{License: ...}
\NormalTok{Encoding: UTF{-}8}
\end{Highlighting}
\end{Shaded}

On y retrouvera notamment :

\begin{itemize}
\tightlist
\item
  le nom de la librairie;
\item
  son titre;
\item
  sa version;
\item
  ses auteurs;
\item
  sa description;
\item
  sa licence;
\item
  ses dépendances.
\end{itemize}

Nous pourrions par exemple remplacer les informations génériques par :

\begin{Shaded}
\begin{Highlighting}[]
\NormalTok{Package: mat8186demo}
\NormalTok{Title: Demonstration for MAT8186}
\NormalTok{Version: 0.0.0.9000}
\NormalTok{Authors@R:}
\NormalTok{    person("Cedric", "Beaulac",}
\NormalTok{           email = "beaulac.cedric@uqam.ca",}
\NormalTok{           role = c("aut", "cre"))}
\NormalTok{Description: Small functions used to demonstrate the development}
\NormalTok{    of an R library in MAT8186.}
\NormalTok{Encoding: UTF{-}8}
\end{Highlighting}
\end{Shaded}

Nous reviendrons plus tard sur la structure exacte du fichier
\texttt{DESCRIPTION}, concentrons sur un élément d'intérêt de ce
fichier, la licence.

\section{Ajouter une licence}\label{ajouter-une-licence}

Si nous distribuons notre librairie, nous devons également indiquer
\textbf{ce que les autres utilisateurs ont le droit de faire avec notre
code}.

Par exemple, nous pouvons choisir une licence MIT :

\begin{Shaded}
\begin{Highlighting}[]
\NormalTok{usethis}\SpecialCharTok{::}\FunctionTok{use\_mit\_license}\NormalTok{()}
\end{Highlighting}
\end{Shaded}

\texttt{usethis} crée et modifie automatiquement les fichiers
nécessaires.

Les licences sont expliqués en plus grand détail au
\href{https://r-pkgs.org/license.html}{chapitre 12 du manuel de
référence}, nous n'en reparlons pas vraiment dans ce cours.

\section{Documenter notre fonction}\label{documenter-notre-fonction}

Notre fonction fonctionne mais un utilisateur qui découvre notre
librairie doit également pouvoir comprendre :

\begin{itemize}
\tightlist
\item
  ce que fait la fonction;
\item
  quels arguments elle prend;
\item
  ce qu'elle retourne;
\item
  comment l'utiliser.
\end{itemize}

Souvenez-vous que lorsque vous approchez une nouvelle libraire, nous
utilisons la documentation R \texttt{?Library::function} pour mieux
comprendre son fonctionnement.

Quand nous mettons au point une libraire, c'est à nous d'écrire la
documentation.

Nous allons utiliser \textbf{\texttt{roxygen2}} pour documenter notre
fonction directement dans son fichier R.

Il faut ouvrir le fichier R en question
\texttt{\textasciitilde{}\textbackslash{}mat8186\textbackslash{}CR.R},
cliquer l'onglet \emph{Code} puis \emph{insert roxygen skeleton}. Ça va
crée une entête vide, que l'on pourra remplir.

\begin{Shaded}
\begin{Highlighting}[]
\CommentTok{\#\textquotesingle{} Centre{-}réduire un vecteur}
\CommentTok{\#\textquotesingle{}}
\CommentTok{\#\textquotesingle{} Transforme linéairement un vecteur numérique afin qu\textquotesingle{}il soit de moyenne 0}
\CommentTok{\#\textquotesingle{} et d\textquotesingle{}écart{-}type 1.}
\CommentTok{\#\textquotesingle{}}
\CommentTok{\#\textquotesingle{} @param x Un vecteur numérique.}
\CommentTok{\#\textquotesingle{}}
\CommentTok{\#\textquotesingle{} @return Un vecteur numérique de même longueur que \textasciigrave{}x\textasciigrave{}.}
\CommentTok{\#\textquotesingle{}}
\CommentTok{\#\textquotesingle{} @examples}
\CommentTok{\#\textquotesingle{} CR(c(10, 20, 30))}
\CommentTok{\#\textquotesingle{}}
\CommentTok{\#\textquotesingle{} @export}
\NormalTok{CR }\OtherTok{\textless{}{-}} \ControlFlowTok{function}\NormalTok{(x) \{}
\NormalTok{  (x }\SpecialCharTok{{-}} \FunctionTok{mean}\NormalTok{(x)) }\SpecialCharTok{/} \FunctionTok{sd}\NormalTok{(x)}
\NormalTok{\}}
\end{Highlighting}
\end{Shaded}

Nous pouvons ensuite générer la documentation grâce à la libraire
\texttt{devtools} :

\begin{Shaded}
\begin{Highlighting}[]
\NormalTok{devtools}\SpecialCharTok{::}\FunctionTok{document}\NormalTok{()}
\end{Highlighting}
\end{Shaded}

Cette commande crée notamment un répertoire :

\begin{Shaded}
\begin{Highlighting}[]
\NormalTok{man/}
\end{Highlighting}
\end{Shaded}

contenant les pages d'aide de notre librairie.

Bon, on se rapelle notre boucle de développement, allons-y:

\begin{Shaded}
\begin{Highlighting}[]
\NormalTok{modifier le code}
\NormalTok{       ↓}
\NormalTok{   load\_all()}
\NormalTok{       ↓}
\NormalTok{tester la fonction}
\NormalTok{       ↓}
\NormalTok{modifier le code}
\NormalTok{       ↓}
\NormalTok{   load\_all()}
\NormalTok{       ↓}
\NormalTok{      ...}
\end{Highlighting}
\end{Shaded}

\begin{Shaded}
\begin{Highlighting}[]
\NormalTok{devtools}\SpecialCharTok{::}\FunctionTok{load\_all}\NormalTok{()}
\NormalTok{?CR}
\end{Highlighting}
\end{Shaded}

Notre fonction possède maintenant une documentation!

Il s'agit peut-être de bon moment pour utiliser Git et laisser une trace
des ces premières modifications.

\section*{Git}\label{git}
\addcontentsline{toc}{section}{Git}

\markright{Git}

On va \emph{staged} les documents modifiés que nous désirons sauvegarder
et étiquette comme provenant de la même ronde de modification.

Puis quand ces documents furent bien choisies, nous allons \emph{commit}
les changement et les décrire à l'aide de commentaires.

Finalement, lorsque l'on \emph{push} ces modifications, la sauvegarde
est déployé!

\section{Ajouter un test}\label{ajouter-un-test}

Pour l'instant, nous avons testé notre fonction manuellement.

Mais comment savoir si une modification future brise accidentellement
son comportement?

Nous allons ajouter des \textbf{tests automatisés}.

\begin{Shaded}
\begin{Highlighting}[]
\NormalTok{usethis}\SpecialCharTok{::}\FunctionTok{use\_testthat}\NormalTok{()}
\end{Highlighting}
\end{Shaded}

Puis créons un fichier de tests associé à notre fonction :

\begin{Shaded}
\begin{Highlighting}[]
\NormalTok{usethis}\SpecialCharTok{::}\FunctionTok{use\_test}\NormalTok{(}\StringTok{"CR"}\NormalTok{)}
\end{Highlighting}
\end{Shaded}

Nous pouvons y écrire :

\begin{Shaded}
\begin{Highlighting}[]
\FunctionTok{test\_that}\NormalTok{(}\StringTok{"CR transforme correctement un vecteur"}\NormalTok{, \{}

  \FunctionTok{expect\_equal}\NormalTok{(}
    \FunctionTok{CR}\NormalTok{(}\FunctionTok{c}\NormalTok{(}\DecValTok{10}\NormalTok{, }\DecValTok{20}\NormalTok{, }\DecValTok{30}\NormalTok{)),}
    \FunctionTok{c}\NormalTok{(}\SpecialCharTok{{-}}\DecValTok{1}\NormalTok{, }\DecValTok{0}\NormalTok{, }\DecValTok{1}\NormalTok{)}
\NormalTok{  )}

\NormalTok{\})}
\end{Highlighting}
\end{Shaded}

Nous lançons ensuite les tests avec :

\begin{Shaded}
\begin{Highlighting}[]
\NormalTok{devtools}\SpecialCharTok{::}\FunctionTok{test}\NormalTok{()}
\end{Highlighting}
\end{Shaded}

Si notre fonction est modifiée plus tard et ne produit plus le résultat
attendu, le test nous avertira.

\section{Utiliser une autre
librairie}\label{utiliser-une-autre-librairie}

Une librairie R n'a évidemment pas besoin de tout réimplémenter
elle-même.

Elle peut utiliser des fonctions provenant d'autres librairies, on
appele ça des dépendances. Ici, c'est une des grandes forces des
librairies R (en contraste avec Python), la gestion des dépendances.

Supposons par exemple que nous souhaitions utiliser la librairie
\texttt{scales}.

Nous pouvons déclarer cette dépendance avec :

\begin{Shaded}
\begin{Highlighting}[]
\NormalTok{usethis}\SpecialCharTok{::}\FunctionTok{use\_package}\NormalTok{(}\StringTok{"plyr"}\NormalTok{)}
\end{Highlighting}
\end{Shaded}

Cette commande ajoute notamment \texttt{plyr} au fichier
\texttt{DESCRIPTION}.

\begin{tcolorbox}[enhanced jigsaw, arc=.35mm, title=\textcolor{quarto-callout-tip-color}{\faLightbulb}\hspace{0.5em}{Le brio des dépendances R}, colback=white, left=2mm, rightrule=.15mm, toptitle=1mm, opacityback=0, bottomtitle=1mm, breakable, colframe=quarto-callout-tip-color-frame, coltitle=black, toprule=.15mm, titlerule=0mm, opacitybacktitle=0.6, bottomrule=.15mm, colbacktitle=quarto-callout-tip-color!10!white, leftrule=.75mm]

Nous pouvons spécifier une version spécifique de la dépendance et R peut
gérer plusieurs versions de la même librairie pour satisfaire les
dépendances si nécessaire.

\end{tcolorbox}

Nous pourrions ensuite écrire une fonction qui centre-réduit chaque
colonne d'une base de données comme :

\begin{Shaded}
\begin{Highlighting}[]
\NormalTok{CR}\SpecialCharTok{{-}}\NormalTok{Donnee }\OtherTok{\textless{}{-}} \ControlFlowTok{function}\NormalTok{(X) \{}
  \FunctionTok{apply}\NormalTok{(x, }\DecValTok{2}\NormalTok{,CR)}
\NormalTok{\}}

\NormalTok{usethis}\SpecialCharTok{::}\FunctionTok{use\_r}\NormalTok{(}\StringTok{"CR{-}Donnee"}\NormalTok{)}
\end{Highlighting}
\end{Shaded}

Et maintenant?

Nous relançons notre cycle :

\begin{Shaded}
\begin{Highlighting}[]
\NormalTok{devtools}\SpecialCharTok{::}\FunctionTok{load\_all}\NormalTok{()}

\NormalTok{devtools}\SpecialCharTok{::}\FunctionTok{test}\NormalTok{()}

\NormalTok{devtools}\SpecialCharTok{::}\FunctionTok{check}\NormalTok{()}
\end{Highlighting}
\end{Shaded}

Nos tests nous permettent de vérifier que le comportement de notre
fonction n'a pas changé, même si son implémentation interne est
maintenant différente.

\section{Installer notre librairie}\label{installer-notre-librairie}

Jusqu'à maintenant, \texttt{load\_all()} nous permettait de simuler le
chargement de notre librairie pendant son développement.

Nous pouvons maintenant réellement l'installer :

\begin{Shaded}
\begin{Highlighting}[]
\NormalTok{devtools}\SpecialCharTok{::}\FunctionTok{install}\NormalTok{()}
\end{Highlighting}
\end{Shaded}

Après avoir redémarré R, nous pouvons l'utiliser comme n'importe quelle
autre librairie :

\begin{Shaded}
\begin{Highlighting}[]
\FunctionTok{library}\NormalTok{(mat8186demo)}

\FunctionTok{CR}\NormalTok{(}\FunctionTok{c}\NormalTok{(}\DecValTok{5}\NormalTok{, }\DecValTok{10}\NormalTok{, }\DecValTok{15}\NormalTok{))}
\end{Highlighting}
\end{Shaded}

C'est une étape importante.

Nous ne possédons plus simplement \textbf{un fichier \texttt{.R}
contenant une fonction}.

Nous avons une \textbf{librairie R installable}.

\section*{Une dernière
vérification}\label{une-derniuxe8re-vuxe9rification}
\addcontentsline{toc}{section}{Une dernière vérification}

\markright{Une dernière vérification}

Nous avons maintenant ajouté plusieurs composantes à notre librairie.

Faisons donc une nouvelle vérification :

\begin{Shaded}
\begin{Highlighting}[]
\NormalTok{devtools}\SpecialCharTok{::}\FunctionTok{check}\NormalTok{()}
\end{Highlighting}
\end{Shaded}

Nous avons donc déjà :

\begin{itemize}
\tightlist
\item
  une librairie R;
\item
  deux fonctions;
\item
  de la documentation;
\item
  une licence;
\item
  des tests;
\item
  une dépendance;
\item
  un historique Git.
\end{itemize}

\section{Et GitHub?}\label{et-github}

Notre dépôt Git existe actuellement \textbf{sur notre ordinateur}.

Une prochaine étape consistera à créer une copie distante de ce dépôt
sur \textbf{GitHub}.

Conceptuellement :

\begin{Shaded}
\begin{Highlighting}[]
\NormalTok{Notre ordinateur}
\NormalTok{     │}
\NormalTok{     │  Git}
\NormalTok{     ▼}
\NormalTok{Historique local}
\NormalTok{     │}
\NormalTok{     │  push / pull}
\NormalTok{     ▼}
\NormalTok{    GitHub}
\end{Highlighting}
\end{Shaded}

Nous verrons dans quelques instants comment mettre cela en place.

\begin{tcolorbox}[enhanced jigsaw, arc=.35mm, title=\textcolor{quarto-callout-warning-color}{\faExclamationTriangle}\hspace{0.5em}{Git n'est pas GitHub}, colback=white, left=2mm, rightrule=.15mm, toptitle=1mm, opacityback=0, bottomtitle=1mm, breakable, colframe=quarto-callout-warning-color-frame, coltitle=black, toprule=.15mm, titlerule=0mm, opacitybacktitle=0.6, bottomrule=.15mm, colbacktitle=quarto-callout-warning-color!10!white, leftrule=.75mm]

\textbf{Git} est le système qui permet de conserver l'historique des
versions de notre projet.

\textbf{GitHub} est un service en ligne qui permet notamment d'héberger
des dépôts Git et de collaborer avec d'autres personnes.

On peut utiliser Git sans GitHub.

\end{tcolorbox}

\section{Le cycle de développement}\label{le-cycle-de-duxe9veloppement}

Nous venons de voir énormément de choses très rapidement.

Le but n'était pas de mémoriser toutes les commandes.

Le message principal est de comprendre le \textbf{cycle général de
développement} d'une librairie :

\begin{Shaded}
\begin{Highlighting}[]
\NormalTok{Créer la librairie}
\NormalTok{        ↓}
\NormalTok{Écrire une fonction}
\NormalTok{        ↓}
\NormalTok{    load\_all()}
\NormalTok{        ↓}
\NormalTok{Tester manuellement}
\NormalTok{        ↓}
\NormalTok{Documenter}
\NormalTok{        ↓}
\NormalTok{Tester automatiquement}
\NormalTok{        ↓}
\NormalTok{     check()}
\NormalTok{        ↓}
\NormalTok{     commit}
\NormalTok{        ↓}
\NormalTok{Modifier le code}
\NormalTok{        ↓}
\NormalTok{     recommencer}
\end{Highlighting}
\end{Shaded}

Une librairie n'est donc pas simplement \textbf{la forme finale que l'on
donne à notre code lorsqu'il est terminé}.

La librairie devient elle-même \textbf{l'environnement dans lequel nous
développons notre code}.

\begin{tcolorbox}[enhanced jigsaw, arc=.35mm, title=\textcolor{quarto-callout-tip-color}{\faLightbulb}\hspace{0.5em}{Pour résumer}, colback=white, left=2mm, rightrule=.15mm, toptitle=1mm, opacityback=0, bottomtitle=1mm, breakable, colframe=quarto-callout-tip-color-frame, coltitle=black, toprule=.15mm, titlerule=0mm, opacitybacktitle=0.6, bottomrule=.15mm, colbacktitle=quarto-callout-tip-color!10!white, leftrule=.75mm]

Nous venons essentiellement de faire le cours au complet en accéléré.

Maintenant, nous allons revenir en arrière et apprendre à faire chacune
de ces étapes correctement.

\end{tcolorbox}

\section*{Et maintenant?}\label{et-maintenant}
\addcontentsline{toc}{section}{Et maintenant?}

\markright{Et maintenant?}

Avant de poursuivre le développement de notre librairie, nous allons
revenir sur deux concepts que nous venons déjà d'utiliser sans vraiment
les expliquer :

\begin{enumerate}
\def\labelenumi{\arabic{enumi}.}
\tightlist
\item
  les \textbf{projets RStudio et l'organisation des fichiers};
\item
  \textbf{Git et GitHub}.
\end{enumerate}

Nous pourrons ensuite revenir au développement de notre librairie avec
une meilleure compréhension de notre environnement de travail.

\bookmarksetup{startatroot}

\chapter{Projet R}\label{projet-r}

Avant d'introduire GitHub et la publication en ligne de notre librairie,
discutons rapidement d'un outil qui nous sera très utile pour organiser
notre travail : le \textbf{projet RStudio}.

\section{Le problème}\label{le-probluxe8me}

On commence souvent un projet en définissant manuellement notre
répertoire de travail :

\begin{Shaded}
\begin{Highlighting}[]
\FunctionTok{setwd}\NormalTok{(}\StringTok{"C:/Users/MYPC/Desktop/MAT8186"}\NormalTok{)}
\end{Highlighting}
\end{Shaded}

Ça fonctionne très bien\ldots{} sur mon ordinateur!

Le problème apparaît rapidement lorsque :

\begin{itemize}
\tightlist
\item
  le projet est déplacé vers un autre répertoire;
\item
  le projet est synchronisé avec un service infonuagique;
\item
  le projet est partagé avec un collègue;
\item
  le projet est téléchargé sur un autre ordinateur.
\end{itemize}

Par exemple, le chemin

\begin{Shaded}
\begin{Highlighting}[]
\NormalTok{C:/Users/MYPC/Desktop/MAT8186}
\end{Highlighting}
\end{Shaded}

n'existe probablement pas sur votre ordinateur.

On aimerait donc que notre code puisse fonctionner
\textbf{indépendamment de l'endroit où le projet est enregistré}.

\begin{tcolorbox}[enhanced jigsaw, arc=.35mm, title=\textcolor{quarto-callout-important-color}{\faExclamation}\hspace{0.5em}{Un principe important}, colback=white, left=2mm, rightrule=.15mm, toptitle=1mm, opacityback=0, bottomtitle=1mm, breakable, colframe=quarto-callout-important-color-frame, coltitle=black, toprule=.15mm, titlerule=0mm, opacitybacktitle=0.6, bottomrule=.15mm, colbacktitle=quarto-callout-important-color!10!white, leftrule=.75mm]

Notre code ne devrait pas dépendre de l'emplacement exact du projet sur
notre ordinateur.

\end{tcolorbox}

\section{Le projet RStudio}\label{le-projet-rstudio}

Un \textbf{projet RStudio} permet de regrouper dans un même répertoire
tout ce qui est associé à une tâche :

\begin{Shaded}
\begin{Highlighting}[]
\NormalTok{mon\_projet/}
\NormalTok{├── mon\_projet.Rproj}
\NormalTok{├── R/}
\NormalTok{├── data/}
\NormalTok{├── figures/}
\NormalTok{└── notes.qmd}
\end{Highlighting}
\end{Shaded}

Le fichier \texttt{.Rproj} identifie ce répertoire comme un projet
RStudio.

Lorsque nous ouvrons ce fichier, RStudio ouvre le projet et utilise
automatiquement \textbf{la racine du projet comme répertoire de
travail}.

Autrement dit, si notre projet est :

\begin{Shaded}
\begin{Highlighting}[]
\NormalTok{C:/Users/MYPC/Desktop/MAT8186/}
\end{Highlighting}
\end{Shaded}

alors :

\begin{Shaded}
\begin{Highlighting}[]
\FunctionTok{getwd}\NormalTok{()}
\end{Highlighting}
\end{Shaded}

retournera quelque chose comme :

\begin{Shaded}
\begin{Highlighting}[]
\NormalTok{C:/Users/MYPC/Desktop/MAT8186}
\end{Highlighting}
\end{Shaded}

Mais si un autre utilisateur place exactement le même projet dans :

\begin{Shaded}
\begin{Highlighting}[]
\NormalTok{C:/Users/YOURPC/Documents/MAT8186/}
\end{Highlighting}
\end{Shaded}

la racine de son projet sera automatiquement :

\begin{Shaded}
\begin{Highlighting}[]
\NormalTok{C:/Users/YOURPC/Documents/MAT8186}
\end{Highlighting}
\end{Shaded}

Il n'est donc plus nécessaire d'écrire un \texttt{setwd()} spécifique à
chaque ordinateur.

\section*{Les chemins relatifs}\label{les-chemins-relatifs}
\addcontentsline{toc}{section}{Les chemins relatifs}

\markright{Les chemins relatifs}

C'est ici qu'apparaît l'un des principaux avantages du projet.

Supposons que notre structure soit :

\begin{Shaded}
\begin{Highlighting}[]
\NormalTok{MAT8186/}
\NormalTok{├── MAT8186.Rproj}
\NormalTok{├── data/}
\NormalTok{│   └── donnees.csv}
\NormalTok{├── R/}
\NormalTok{│   └── analyse.R}
\NormalTok{└── notes.qmd}
\end{Highlighting}
\end{Shaded}

Plutôt que d'écrire :

\begin{Shaded}
\begin{Highlighting}[]
\FunctionTok{read.csv}\NormalTok{(}
  \StringTok{"C:/Users/MYPC/Desktop/MAT8186/data/donnees.csv"}
\NormalTok{)}
\end{Highlighting}
\end{Shaded}

nous pouvons simplement écrire :

\begin{Shaded}
\begin{Highlighting}[]
\FunctionTok{read.csv}\NormalTok{(}\StringTok{"data/donnees.csv"}\NormalTok{)}
\end{Highlighting}
\end{Shaded}

Le chemin est maintenant \textbf{relatif à la racine du projet}.

Si nous déplaçons tout le projet ailleurs, le chemin

\begin{Shaded}
\begin{Highlighting}[]
\NormalTok{data/donnees.csv}
\end{Highlighting}
\end{Shaded}

demeure valide.

Il fonctionnera également sur l'ordinateur d'un collègue qui possède la
même structure de projet.

\begin{tcolorbox}[enhanced jigsaw, arc=.35mm, title=\textcolor{quarto-callout-tip-color}{\faLightbulb}\hspace{0.5em}{Une bonne règle}, colback=white, left=2mm, rightrule=.15mm, toptitle=1mm, opacityback=0, bottomtitle=1mm, breakable, colframe=quarto-callout-tip-color-frame, coltitle=black, toprule=.15mm, titlerule=0mm, opacitybacktitle=0.6, bottomrule=.15mm, colbacktitle=quarto-callout-tip-color!10!white, leftrule=.75mm]

Si votre code contient :

\begin{Shaded}
\begin{Highlighting}[]
\FunctionTok{setwd}\NormalTok{(}\StringTok{"C:/Users/mon\_nom/..."}\NormalTok{)}
\end{Highlighting}
\end{Shaded}

ou un autre chemin propre à votre ordinateur, demandez-vous s'il ne
serait pas préférable de travailler à partir d'un projet et d'utiliser
un chemin relatif.

\end{tcolorbox}

\section{Partir à neuf}\label{partir-uxe0-neuf}

R et RStudio peuvent conserver automatiquement une partie de notre
environnement de travail entre deux sessions.

Deux fichiers que vous avez peut-être déjà rencontrés sont :

\begin{Shaded}
\begin{Highlighting}[]
\NormalTok{.RData}
\NormalTok{.Rhistory}
\end{Highlighting}
\end{Shaded}

\texttt{.RData} peut contenir les objets présents dans notre
environnement R à la fin d'une session, tandis que \texttt{.Rhistory}
peut conserver l'historique des commandes exécutées.

Cela peut sembler pratique.

Cependant, dans un projet reproductible, nous voulons idéalement être
capables de reconstruire notre environnement \textbf{à partir de notre
code}.

Par exemple, si notre analyse nécessite l'objet \texttt{donnees}, nous
préférons avoir :

\begin{Shaded}
\begin{Highlighting}[]
\NormalTok{donnees }\OtherTok{\textless{}{-}} \FunctionTok{read.csv}\NormalTok{(}\StringTok{"data/donnees.csv"}\NormalTok{)}
\end{Highlighting}
\end{Shaded}

dans notre code plutôt que de dépendre d'un objet \texttt{donnees} qui
serait mystérieusement réapparu depuis une ancienne session R.

L'idée est donc de travailler \textbf{avec intention} : les objets
nécessaires à notre analyse devraient pouvoir être recréés à partir des
fichiers et du code contenus dans notre projet.

Pour cette raison, il est généralement recommandé de ne pas restaurer
automatiquement \texttt{.RData} au démarrage et de ne pas sauvegarder
automatiquement l'environnement de travail à la fermeture de RStudio.

Nous pouvons configurer ce comportement dans :

\begin{Shaded}
\begin{Highlighting}[]
\NormalTok{Tools {-}\textgreater{} Global Options {-}\textgreater{} General}
\end{Highlighting}
\end{Shaded}

et choisir notamment :

\begin{Shaded}
\begin{Highlighting}[]
\NormalTok{Restore .RData into workspace at startup : NON}

\NormalTok{Save workspace to .RData on exit : Never}
\end{Highlighting}
\end{Shaded}

\begin{tcolorbox}[enhanced jigsaw, arc=.35mm, title=\textcolor{quarto-callout-note-color}{\faInfo}\hspace{0.5em}{Reproductibilité}, colback=white, left=2mm, rightrule=.15mm, toptitle=1mm, opacityback=0, bottomtitle=1mm, breakable, colframe=quarto-callout-note-color-frame, coltitle=black, toprule=.15mm, titlerule=0mm, opacitybacktitle=0.6, bottomrule=.15mm, colbacktitle=quarto-callout-note-color!10!white, leftrule=.75mm]

Une bonne question à se poser est :

\begin{quote}
Si je ferme complètement RStudio, que je rouvre mon projet et que je
réexécute mon code, est-ce que je peux reproduire mes résultats?
\end{quote}

Si la réponse est oui, nous sommes sur la bonne voie.

\end{tcolorbox}

\section{Créer un projet}\label{cruxe9er-un-projet}

Lorsque nous commençons sérieusement une nouvelle tâche -- une analyse
de données, la rédaction de notes Quarto ou le développement d'une
librairie -- nous pouvons créer un projet avec :

\begin{Shaded}
\begin{Highlighting}[]
\NormalTok{File {-}\textgreater{} New Project}
\end{Highlighting}
\end{Shaded}

RStudio nous offre notamment trois possibilités :

\begin{Shaded}
\begin{Highlighting}[]
\NormalTok{New Directory}
\NormalTok{Existing Directory}
\NormalTok{Version Control}
\end{Highlighting}
\end{Shaded}

Nous pouvons donc :

\begin{itemize}
\tightlist
\item
  créer simultanément un nouveau répertoire et un nouveau projet;
\item
  transformer un répertoire existant en projet RStudio;
\item
  récupérer un projet existant à partir d'un dépôt Git.
\end{itemize}

Lors de la création du projet, RStudio crée notamment un fichier :

\begin{Shaded}
\begin{Highlighting}[]
\NormalTok{nom\_du\_projet.Rproj}
\end{Highlighting}
\end{Shaded}

Par la suite, lorsque nous voulons travailler sur ce projet, nous
ouvrons \textbf{le projet lui-même} plutôt que d'ouvrir RStudio puis de
chercher nos fichiers.

Nous obtenons ainsi un environnement de travail bien délimité :

\begin{Shaded}
\begin{Highlighting}[]
\NormalTok{Projet A}
\NormalTok{├── données}
\NormalTok{├── code}
\NormalTok{├── résultats}
\NormalTok{└── environnement de travail}

\NormalTok{Projet B}
\NormalTok{├── données}
\NormalTok{├── code}
\NormalTok{├── résultats}
\NormalTok{└── environnement de travail}
\end{Highlighting}
\end{Shaded}

Les différents projets demeurent ainsi bien séparés les uns des autres.

\section*{Et notre librairie?}\label{et-notre-librairie}
\addcontentsline{toc}{section}{Et notre librairie?}

\markright{Et notre librairie?}

Nous avons en fait déjà utilisé un projet RStudio!

Lorsque nous avons créé notre librairie avec :

\begin{Shaded}
\begin{Highlighting}[]
\NormalTok{usethis}\SpecialCharTok{::}\FunctionTok{create\_package}\NormalTok{(}\StringTok{"mat8186demo"}\NormalTok{)}
\end{Highlighting}
\end{Shaded}

nous avons obtenu quelque chose comme :

\begin{Shaded}
\begin{Highlighting}[]
\NormalTok{mat8186demo/}
\NormalTok{├── mat8186demo.Rproj}
\NormalTok{├── DESCRIPTION}
\NormalTok{├── NAMESPACE}
\NormalTok{└── R/}
\end{Highlighting}
\end{Shaded}

Notre librairie est donc également un \textbf{projet RStudio}.

Et lors de notre survol du développement d'une librairie, nous avons
ajouté :

\begin{Shaded}
\begin{Highlighting}[]
\NormalTok{usethis}\SpecialCharTok{::}\FunctionTok{use\_git}\NormalTok{()}
\end{Highlighting}
\end{Shaded}

Notre même répertoire est maintenant :

\begin{Shaded}
\begin{Highlighting}[]
\NormalTok{un répertoire}
\NormalTok{      +}
\NormalTok{un projet RStudio}
\NormalTok{      +}
\NormalTok{une librairie R}
\NormalTok{      +}
\NormalTok{un dépôt Git}
\end{Highlighting}
\end{Shaded}

Il nous reste maintenant à comprendre un peu mieux ce dernier élément.

\section{Git et GitHub}\label{git-et-github}

Nous avons déjà créé un dépôt Git sans vraiment expliquer ce que cela
signifie.

Il est maintenant temps de revenir sur \textbf{Git}, puis de comprendre
son lien avec \textbf{GitHub}.

\bookmarksetup{startatroot}

\chapter{Git : garder une trace de notre
travail}\label{git-garder-une-trace-de-notre-travail}

Nous venons de créer une librairie et d'utiliser Git presque
immédiatement avec :

\begin{Shaded}
\begin{Highlighting}[]
\NormalTok{usethis}\SpecialCharTok{::}\FunctionTok{use\_git}\NormalTok{()}
\end{Highlighting}
\end{Shaded}

Mais qu'avons-nous réellement fait?

Git est un \textbf{système de contrôle des versions} (\emph{version
control system}).

Son rôle est de conserver de manière structurée l'évolution d'un
ensemble de fichiers.

Une excellente référence pour Git dans le contexte de R est le livre
\href{https://happygitwithr.com/}{Happy Git and GitHub for the useR} de
Jenny Bryan et collaborateurs. Nous avons conçu ce chapitre en
s'inspirant de cette source.

\begin{tcolorbox}[enhanced jigsaw, arc=.35mm, title=\textcolor{quarto-callout-note-color}{\faInfo}\hspace{0.5em}{Une approche volontairement minimale}, colback=white, left=2mm, rightrule=.15mm, toptitle=1mm, opacityback=0, bottomtitle=1mm, breakable, colframe=quarto-callout-note-color-frame, coltitle=black, toprule=.15mm, titlerule=0mm, opacitybacktitle=0.6, bottomrule=.15mm, colbacktitle=quarto-callout-note-color!10!white, leftrule=.75mm]

Git est un outil extrêmement puissant et possède énormément de
fonctionnalités.

Dans ce cours, notre objectif n'est pas de devenir experts de Git.

Nous voulons apprendre \textbf{le plus petit nombre d'opérations
nécessaires pour intégrer Git naturellement à notre travail de
mise-au-point de librairie}.

\end{tcolorbox}

\section*{Pourquoi utiliser un système de contrôle des
versions?}\label{pourquoi-utiliser-un-systuxe8me-de-contruxf4le-des-versions}
\addcontentsline{toc}{section}{Pourquoi utiliser un système de contrôle
des versions?}

\markright{Pourquoi utiliser un système de contrôle des versions?}

Vous avez probablement déjà utilisé un système de contrôle des versions
artisanal :

\begin{Shaded}
\begin{Highlighting}[]
\NormalTok{analyse.R}
\NormalTok{analyse\_final.R}
\NormalTok{analyse\_final2.R}
\NormalTok{analyse\_final\_VRAIMENT\_FINAL.R}
\NormalTok{analyse\_final\_VRAIMENT\_FINAL\_corrige.R}
\end{Highlighting}
\end{Shaded}

Après quelques semaines, il devient difficile de répondre à des
questions simples :

\begin{itemize}
\tightlist
\item
  Quelle est la version actuelle?
\item
  Qu'est-ce qui a changé entre deux versions?
\item
  Pourquoi avons-nous effectué cette modification?
\item
  Qui a effectué la modification?
\item
  Est-il possible de revenir à une version antérieure?
\item
  Est-ce que cette ancienne version fonctionnait mieux?
\end{itemize}

Git propose une solution structurée à ce problème.

Plutôt que de créer continuellement de nouvelles copies du même fichier,
nous en conservons qu'un seul et nous conservons, par Git, son
évolution.

\section{Un répertoire Git}\label{un-ruxe9pertoire-git}

L'unité fondamentale de Git est le \textbf{répertoire}
(\emph{repository} ou simplement \emph{repo}).

Un répertoire Git est essentiellement un fichier dont Git suit
l'évolution.

Par exemple :

\begin{Shaded}
\begin{Highlighting}[]
\NormalTok{mat8186demo/}
\NormalTok{├── DESCRIPTION}
\NormalTok{├── NAMESPACE}
\NormalTok{├── R/}
\NormalTok{│   └── rescale01.R}
\NormalTok{├── tests/}
\NormalTok{├── man/}
\NormalTok{└── .git/}
\end{Highlighting}
\end{Shaded}

Le répertoire caché :

\begin{Shaded}
\begin{Highlighting}[]
\NormalTok{.git/}
\end{Highlighting}
\end{Shaded}

contient l'information dont Git a besoin pour conserver l'historique du
projet.

\begin{tcolorbox}[enhanced jigsaw, arc=.35mm, title=\textcolor{quarto-callout-tip-color}{\faLightbulb}\hspace{0.5em}{Un dépôt reste un répertoire}, colback=white, left=2mm, rightrule=.15mm, toptitle=1mm, opacityback=0, bottomtitle=1mm, breakable, colframe=quarto-callout-tip-color-frame, coltitle=black, toprule=.15mm, titlerule=0mm, opacitybacktitle=0.6, bottomrule=.15mm, colbacktitle=quarto-callout-tip-color!10!white, leftrule=.75mm]

Transformer un répertoire en dépôt Git ne change pas fondamentalement
vos fichiers.

Vous pouvez toujours ouvrir, modifier, déplacer et utiliser vos fichiers
normalement.

Git ajoute simplement une couche permettant d'en conserver l'historique.

\end{tcolorbox}

Dans notre cas :

\begin{Shaded}
\begin{Highlighting}[]
\NormalTok{mat8186demo/}
\end{Highlighting}
\end{Shaded}

est donc simultanément :

\begin{Shaded}
\begin{Highlighting}[]
\NormalTok{un répertoire}
\NormalTok{      +}
\NormalTok{un projet RStudio}
\NormalTok{      +}
\NormalTok{une librairie R}
\NormalTok{      +}
\NormalTok{un dépôt Git}
\end{Highlighting}
\end{Shaded}

\section{Le commit}\label{le-commit}

Git ne conserve pas automatiquement chaque caractère que vous tapez.

Vous décidez plutôt quand une étape de votre travail mérite d'être
enregistrée.

Cet enregistrement s'appelle un \textbf{commit}.

On peut voir un commit comme une photographie de l'état du projet à un
moment jugé important.

Par exemple :

\begin{Shaded}
\begin{Highlighting}[]
\NormalTok{A    Initial package structure}

\NormalTok{B    Add CR function}

\NormalTok{C    Add documentation}

\NormalTok{D    Add tests for CR}
\end{Highlighting}
\end{Shaded}

Chaque commit représente donc une étape identifiable du développement de
notre librairie.

\section*{Pourquoi ne pas simplement
sauvegarder?}\label{pourquoi-ne-pas-simplement-sauvegarder}
\addcontentsline{toc}{section}{Pourquoi ne pas simplement sauvegarder?}

\markright{Pourquoi ne pas simplement sauvegarder?}

Sauvegarder un fichier et créer un commit sont deux choses différentes.

Lorsque nous sauvegardons un document (\texttt{Ctrl\ +\ S}), nous
remplaçons sa version précédente.

Un commit sauvegarde la version actuelle du fichier tout en conservant
la liste de toutes les modifications faites sur ce fichier. Si nos
fichiers sont constamment sauvegarder par un \textbf{commit}, nous
pouvons donc observer la liste complète des modifications faites sur ce
document.

\section*{Un commit concerne le
projet}\label{un-commit-concerne-le-projet}
\addcontentsline{toc}{section}{Un commit concerne le projet}

\markright{Un commit concerne le projet}

Un aspect important de Git est qu'un commit peut contenir des
modifications apportées à \textbf{plusieurs fichiers simultanément}.

Supposons que nous modifions :

\begin{Shaded}
\begin{Highlighting}[]
\NormalTok{R/CR.R}
\NormalTok{tests/testthat/CR.R}
\NormalTok{DESCRIPTION}
\end{Highlighting}
\end{Shaded}

Ces trois changements peuvent faire partie du même commit si, ensemble,
ils représentent une modification logique du projet.

Par exemple :

\begin{Shaded}
\begin{Highlighting}[]
\NormalTok{Improve CR and update DESCRIPTION}
\end{Highlighting}
\end{Shaded}

\section{Les trois états qui nous
intéressent}\label{les-trois-uxe9tats-qui-nous-intuxe9ressent}

Pour commencer, nous pouvons utiliser un modèle mental très simple.

Un fichier peut être :

\begin{Shaded}
\begin{Highlighting}[]
\NormalTok{modifié}
\NormalTok{   ↓}
\NormalTok{staged}
\NormalTok{   ↓}
\NormalTok{committed}
\end{Highlighting}
\end{Shaded}

\section*{Étape 1: Modifié}\label{uxe9tape-1-modifiuxe9}
\addcontentsline{toc}{section}{Étape 1: Modifié}

\markright{Étape 1: Modifié}

Nous changeons un fichier normalement dans RStudio.

Git remarque :

\begin{quote}
Ce fichier est différent de la dernière version enregistrée.
\end{quote}

Mais rien n'a encore été ajouté à l'historique.

\section*{Étape 2: Staged}\label{uxe9tape-2-staged}
\addcontentsline{toc}{section}{Étape 2: Staged}

\markright{Étape 2: Staged}

Avant de faire un commit, nous choisissons les changements que nous
voulons inclure.

Cette étape s'appelle \textbf{staging}.

Dans l'onglet Git de RStudio, cela correspond essentiellement à cocher
la case \textbf{Staged} devant un fichier.

Nous disons alors à Git :

\begin{quote}
Je veux que cette modification fasse partie de mon prochain commit.
\end{quote}

\section*{Étape 3: Committed}\label{uxe9tape-3-committed}
\addcontentsline{toc}{section}{Étape 3: Committed}

\markright{Étape 3: Committed}

Lorsque nous créons le commit, les changements sélectionnés sont ajoutés
à l'historique du projet.

Nous devons également donner un message au commit.

Par exemple :

\begin{Shaded}
\begin{Highlighting}[]
\NormalTok{Add rescale01 function}
\end{Highlighting}
\end{Shaded}

Notre cycle devient donc :

\begin{Shaded}
\begin{Highlighting}[]
\NormalTok{Modifier}
\NormalTok{   ↓}
\NormalTok{Sauvegarder}
\NormalTok{   ↓}
\NormalTok{Examiner les changements}
\NormalTok{   ↓}
\NormalTok{Stage}
\NormalTok{   ↓}
\NormalTok{Commit}
\end{Highlighting}
\end{Shaded}

\section{Le message de commit}\label{le-message-de-commit}

Chaque commit doit avoir un court message.

Un bon message explique principalement \textbf{l'objectif de la
modification}.

C'est plus difficile que ça en a l'aire\ldots.

\begin{tcolorbox}[enhanced jigsaw, arc=.35mm, title=\textcolor{quarto-callout-tip-color}{\faLightbulb}\hspace{0.5em}{Une bonne habitude}, colback=white, left=2mm, rightrule=.15mm, toptitle=1mm, opacityback=0, bottomtitle=1mm, breakable, colframe=quarto-callout-tip-color-frame, coltitle=black, toprule=.15mm, titlerule=0mm, opacitybacktitle=0.6, bottomrule=.15mm, colbacktitle=quarto-callout-tip-color!10!white, leftrule=.75mm]

Faire plus de plus petit commit\ldots{} (\textbf{petits, fréquents et
cohérents}).

Un commit devrait idéalement représenter une étape logique que vous
pouvez décrire simplement.

\end{tcolorbox}

\section{Git dans RStudio}\label{git-dans-rstudio}

Nous pouvons utiliser Git directement dans le terminal :

\begin{Shaded}
\begin{Highlighting}[]
\FunctionTok{git}\NormalTok{ status}
\FunctionTok{git}\NormalTok{ add}
\FunctionTok{git}\NormalTok{ commit}
\end{Highlighting}
\end{Shaded}

mais RStudio offre une interface graphique pour les opérations que nous
utiliserons le plus souvent.

Lorsqu'un projet RStudio est également un dépôt Git, un onglet
\textbf{Git} apparaît dans RStudio.

Cet onglet nous permet notamment de :

\begin{itemize}
\tightlist
\item
  voir les fichiers modifiés;
\item
  examiner les différences;
\item
  sélectionner les fichiers à inclure dans un commit;
\item
  créer un commit;
\item
  consulter l'historique;
\item
  éventuellement envoyer et recevoir des modifications avec GitHub.
\end{itemize}

\section*{\texorpdfstring{\texttt{git\ status}}{git status}}\label{git-status}
\addcontentsline{toc}{section}{\texttt{git\ status}}

\markright{\texttt{git\ status}}

Même si nous utiliserons principalement l'interface RStudio, une
commande mérite d'être connue dès maintenant :

\begin{Shaded}
\begin{Highlighting}[]
\FunctionTok{git}\NormalTok{ status}
\end{Highlighting}
\end{Shaded}

Elle répond essentiellement à la question :

\begin{quote}
Dans quel état se trouve actuellement mon dépôt?
\end{quote}

On pourrait obtenir :

\begin{Shaded}
\begin{Highlighting}[]
\NormalTok{On branch main}

\NormalTok{Changes not staged for commit:}
\NormalTok{    modified: R/rescale01.R}
\end{Highlighting}
\end{Shaded}

Git nous indique donc qu'un fichier a été modifié depuis le dernier
commit.

Après avoir commité toutes nos modifications :

\begin{Shaded}
\begin{Highlighting}[]
\NormalTok{On branch main}

\NormalTok{nothing to commit, working tree clean}
\end{Highlighting}
\end{Shaded}

L'expression :

\begin{Shaded}
\begin{Highlighting}[]
\NormalTok{working tree clean}
\end{Highlighting}
\end{Shaded}

signifie simplement que notre répertoire de travail ne contient
actuellement aucune modification non enregistrée par Git.

\section{Git n'est pas GitHub}\label{git-nest-pas-github-1}

Il est essentiel de distinguer \textbf{Git} et \textbf{GitHub}.

Git fonctionne entièrement sur votre ordinateur. Nous pourrions utiliser
Git pendant des années sans jamais créer de compte GitHub.

GitHub est essentiellement un service en ligne (\emph{cloud}) qui permet
d'héberger des dépôts Git.

Lorsque nous ajoutons GitHub :

\begin{Shaded}
\begin{Highlighting}[]
\NormalTok{       ordinateur}
\NormalTok{           │}
\NormalTok{           │}
\NormalTok{           ▼}
\NormalTok{      dépôt Git local}
\NormalTok{           │}
\NormalTok{      push │ │ pull}
\NormalTok{           │}
\NormalTok{           ▼}
\NormalTok{        GitHub}
\end{Highlighting}
\end{Shaded}

Git conserve donc l'historique (localement).

GitHub nous permet notamment de :

\begin{itemize}
\tightlist
\item
  conserver une copie distante du dépôt;
\item
  partager notre code;
\item
  collaborer;
\item
  consulter l'historique dans une interface web;
\item
  distribuer notre librairie;
\end{itemize}

\begin{tcolorbox}[enhanced jigsaw, arc=.35mm, title=\textcolor{quarto-callout-important-color}{\faExclamation}\hspace{0.5em}{Deux concepts différents}, colback=white, left=2mm, rightrule=.15mm, toptitle=1mm, opacityback=0, bottomtitle=1mm, breakable, colframe=quarto-callout-important-color-frame, coltitle=black, toprule=.15mm, titlerule=0mm, opacitybacktitle=0.6, bottomrule=.15mm, colbacktitle=quarto-callout-important-color!10!white, leftrule=.75mm]

\textbf{Git = contrôle des versions}

\textbf{GitHub = hébergement et collaboration autour de dépôts Git}

\end{tcolorbox}

\bookmarksetup{startatroot}

\chapter{Le dépôt local et le dépôt
distant}\label{le-duxe9puxf4t-local-et-le-duxe9puxf4t-distant}

Lorsque notre projet existe également sur GitHub, nous avons
essentiellement \textbf{deux copies du dépôt} :

\begin{Shaded}
\begin{Highlighting}[]
\NormalTok{LOCAL                         DISTANT}

\NormalTok{Votre ordinateur              GitHub}
\NormalTok{MAT8186/                       github.com/...}
\NormalTok{    │                              │}
\NormalTok{    └── dépôt Git                  └── dépôt Git}
\end{Highlighting}
\end{Shaded}

Ces deux dépôts peuvent évoluer indépendamment.

Il faut donc parfois transmettre les modifications d'un endroit à
l'autre.

Deux opérations seront particulièrement importantes :

\begin{Shaded}
\begin{Highlighting}[]
\NormalTok{push}
\NormalTok{pull}
\end{Highlighting}
\end{Shaded}

\section{Push}\label{push}

Un \textbf{push} envoie nos nouveaux commits locaux vers GitHub.

\begin{Shaded}
\begin{Highlighting}[]
\NormalTok{ordinateur}
\NormalTok{    │}
\NormalTok{    │ push}
\NormalTok{    ▼}
\NormalTok{ GitHub}
\end{Highlighting}
\end{Shaded}

Un \texttt{push} ne signifie donc pas simplement de sauvegarder un
fichier sur Internet mais bien de transmettre l'historique complet des
modifications de ce fichier sur la version en ligne de ce répertoire
Git.

\section{Pull}\label{pull}

Un \textbf{pull} récupère les nouveaux commits présents sur GitHub.

\begin{Shaded}
\begin{Highlighting}[]
\NormalTok{ordinateur}
\NormalTok{    ▲}
\NormalTok{    │ pull}
\NormalTok{    │}
\NormalTok{ GitHub}
\end{Highlighting}
\end{Shaded}

Cela devient particulièrement important lorsque :

\begin{itemize}
\tightlist
\item
  nous travaillons sur plusieurs ordinateurs;
\item
  plusieurs personnes collaborent sur le même projet;
\item
  des modifications ont été effectuées directement sur GitHub.
\end{itemize}

\section{Notre cycle de travail
quotidien}\label{notre-cycle-de-travail-quotidien}

Pour un projet simple, notre cycle Git/GitHub ressemblera souvent à :

\begin{Shaded}
\begin{Highlighting}[]
\NormalTok{PULL}
\NormalTok{  ↓}
\NormalTok{Modifier des fichiers}
\NormalTok{  ↓}
\NormalTok{Sauvegarder}
\NormalTok{  ↓}
\NormalTok{Examiner les modifications}
\NormalTok{  ↓}
\NormalTok{STAGE}
\NormalTok{  ↓}
\NormalTok{COMMIT}
\NormalTok{  ↓}
\NormalTok{PUSH}
\end{Highlighting}
\end{Shaded}

Puis nous recommençons.

\begin{tcolorbox}[enhanced jigsaw, arc=.35mm, title=\textcolor{quarto-callout-tip-color}{\faLightbulb}\hspace{0.5em}{Une règle simple pour commencer}, colback=white, left=2mm, rightrule=.15mm, toptitle=1mm, opacityback=0, bottomtitle=1mm, breakable, colframe=quarto-callout-tip-color-frame, coltitle=black, toprule=.15mm, titlerule=0mm, opacitybacktitle=0.6, bottomrule=.15mm, colbacktitle=quarto-callout-tip-color!10!white, leftrule=.75mm]

Lorsque vous travaillez avec un dépôt distant :

\textbf{Pull avant de commencer, commit lorsque vous avez accompli une
petite étape cohérente, push régulièrement.}

\end{tcolorbox}

\bookmarksetup{startatroot}

\chapter{Petit exercice}\label{petit-exercice}

Prenons notre librairie \texttt{mat8186demo}.

\subsection{Étape 1}\label{uxe9tape-1}

Ouvrez :

\begin{Shaded}
\begin{Highlighting}[]
\NormalTok{R/CR.R}
\end{Highlighting}
\end{Shaded}

et ajoutez un commentaire :

\begin{Shaded}
\begin{Highlighting}[]
\CommentTok{\# Ma premiere fonction pour centrer{-}réduire!}

\NormalTok{CR }\OtherTok{\textless{}{-}} \ControlFlowTok{function}\NormalTok{(x) \{}
\NormalTok{  (x }\SpecialCharTok{{-}} \FunctionTok{mean}\NormalTok{(x)) }\SpecialCharTok{/} \FunctionTok{sd}\NormalTok{(x)}
\NormalTok{\}}
\end{Highlighting}
\end{Shaded}

Sauvegardez le fichier.

\subsection{Étape 2}\label{uxe9tape-2}

Regardez l'onglet \textbf{Git} de RStudio.

Le fichier devrait apparaître comme modifié.

\subsection{Étape 3}\label{uxe9tape-3}

Examinez le \textbf{Diff}.

Pouvez-vous identifier exactement la ligne ajoutée?

\subsection{Étape 4}\label{uxe9tape-4}

Cochez \textbf{Staged}.

\subsection{Étape 5}\label{uxe9tape-5}

Créez un commit avec le message :

\begin{Shaded}
\begin{Highlighting}[]
\NormalTok{Add comment to CR}
\end{Highlighting}
\end{Shaded}

Nous venons de créer une nouvelle version identifiable de notre projet.

\begin{tcolorbox}[enhanced jigsaw, arc=.35mm, title=\textcolor{quarto-callout-note-color}{\faInfo}\hspace{0.5em}{La suite}, colback=white, left=2mm, rightrule=.15mm, toptitle=1mm, opacityback=0, bottomtitle=1mm, breakable, colframe=quarto-callout-note-color-frame, coltitle=black, toprule=.15mm, titlerule=0mm, opacitybacktitle=0.6, bottomrule=.15mm, colbacktitle=quarto-callout-note-color!10!white, leftrule=.75mm]

Nous savons maintenant ce que Git cherche à accomplir.

La prochaine étape consiste à connecter concrètement :

\textbf{RStudio → Git → GitHub}

et à effectuer nous-mêmes le cycle :

\textbf{modifier → commit → push → pull}.

\end{tcolorbox}

\section{Préparer son
environnement}\label{pruxe9parer-son-environnement}

Avant d'aller plus loin, assurons-nous que tout le monde possède les
outils nécessaires.

Nous aurons besoin de :

\begin{enumerate}
\def\labelenumi{\arabic{enumi}.}
\tightlist
\item
  \textbf{Git} installé sur notre ordinateur;
\item
  un \textbf{compte GitHub};
\item
  une configuration minimale permettant à Git d'identifier l'auteur des
  modifications.
\end{enumerate}

\section*{Installer Git sous Windows}\label{installer-git-sous-windows}
\addcontentsline{toc}{section}{Installer Git sous Windows}

\markright{Installer Git sous Windows}

Sous Windows, nous utiliserons \textbf{Git for Windows}.

\subsection{Étape 1 --- Télécharger
Git}\label{uxe9tape-1-tuxe9luxe9charger-git}

Rendez-vous sur le site officiel de \textbf{Git} et téléchargez la
version de \textbf{Git for Windows} correspondant à votre ordinateur.

Pour la majorité des ordinateurs récents, il s'agira de la version
\textbf{64 bits (x64)}.

\subsection{Étape 2 --- Installer Git}\label{uxe9tape-2-installer-git}

Lancez l'installateur.

Pour ce cours, les \textbf{options proposées par défaut conviennent
généralement}. Il n'est donc pas nécessaire de modifier toutes les
options de l'installateur.

L'installation ajoute notamment :

\begin{itemize}
\tightlist
\item
  Git;
\item
  les commandes Git accessibles dans le terminal;
\item
  \textbf{Git Bash}, un terminal permettant d'utiliser Git sous Windows.
\end{itemize}

\begin{tcolorbox}[enhanced jigsaw, arc=.35mm, title=\textcolor{quarto-callout-tip-color}{\faLightbulb}\hspace{0.5em}{Pas besoin de devenir expert de Git Bash}, colback=white, left=2mm, rightrule=.15mm, toptitle=1mm, opacityback=0, bottomtitle=1mm, breakable, colframe=quarto-callout-tip-color-frame, coltitle=black, toprule=.15mm, titlerule=0mm, opacitybacktitle=0.6, bottomrule=.15mm, colbacktitle=quarto-callout-tip-color!10!white, leftrule=.75mm]

Nous utiliserons principalement Git à travers \textbf{RStudio}.

Le terminal nous servira surtout à vérifier certaines configurations ou
à comprendre ce que RStudio fait derrière son interface.

\end{tcolorbox}

\subsection{Étape 3 --- Vérifier
l'installation}\label{uxe9tape-3-vuxe9rifier-linstallation}

Une fois Git installé, \textbf{redémarrez RStudio}.

Ouvrez ensuite :

\textbf{Tools → Global Options → Git/SVN}

RStudio devrait afficher un chemin dans le champ :

\begin{Shaded}
\begin{Highlighting}[]
\NormalTok{Git executable}
\end{Highlighting}
\end{Shaded}

On peut également ouvrir le terminal de RStudio et exécuter :

\begin{Shaded}
\begin{Highlighting}[]
\FunctionTok{git} \AttributeTok{{-}{-}version}
\end{Highlighting}
\end{Shaded}

Vous devriez obtenir quelque chose comme :

\begin{Shaded}
\begin{Highlighting}[]
\NormalTok{git version 2.x.x}
\end{Highlighting}
\end{Shaded}

La version exacte n'est pas importante.

Si une version est affichée, Git est correctement installé.

\bookmarksetup{startatroot}

\chapter{Créer un compte GitHub}\label{cruxe9er-un-compte-github}

Git fonctionne localement et ne nécessite aucun compte.

Par contre, pour partager nos dépôts et collaborer, nous utiliserons
\textbf{GitHub}.

\section*{Créer le compte}\label{cruxe9er-le-compte}
\addcontentsline{toc}{section}{Créer le compte}

\markright{Créer le compte}

Rendez-vous sur \textbf{GitHub} et créez un compte personnel gratuit.

Vous devrez notamment choisir :

\begin{itemize}
\tightlist
\item
  une adresse courriel;
\item
  un mot de passe;
\item
  un \textbf{nom d'utilisateur GitHub}.
\end{itemize}

\begin{tcolorbox}[enhanced jigsaw, arc=.35mm, title=\textcolor{quarto-callout-important-color}{\faExclamation}\hspace{0.5em}{Choisissez bien votre nom d'utilisateur}, colback=white, left=2mm, rightrule=.15mm, toptitle=1mm, opacityback=0, bottomtitle=1mm, breakable, colframe=quarto-callout-important-color-frame, coltitle=black, toprule=.15mm, titlerule=0mm, opacitybacktitle=0.6, bottomrule=.15mm, colbacktitle=quarto-callout-important-color!10!white, leftrule=.75mm]

Votre nom d'utilisateur GitHub sera visible dans l'adresse de vos
projets et potentiellement dans vos contributions publiques.

Il est donc préférable de choisir un nom relativement professionnel et
que vous serez confortable d'utiliser à long terme.

\end{tcolorbox}

\bookmarksetup{startatroot}

\chapter{Présenter Git à GitHub\ldots{} et à
vous-même}\label{pruxe9senter-git-uxe0-github-et-uxe0-vous-muxeame}

Git doit savoir \textbf{qui effectue les commits}.

Cette information est enregistrée dans la configuration de Git.

Dans le terminal de RStudio, exécutez :

\begin{Shaded}
\begin{Highlighting}[]
\FunctionTok{git}\NormalTok{ config }\AttributeTok{{-}{-}global}\NormalTok{ user.name }\StringTok{"Votre Nom"}
\end{Highlighting}
\end{Shaded}

Par exemple :

\begin{Shaded}
\begin{Highlighting}[]
\FunctionTok{git}\NormalTok{ config }\AttributeTok{{-}{-}global}\NormalTok{ user.name }\StringTok{"Jean Tremblay"}
\end{Highlighting}
\end{Shaded}

Puis indiquez votre adresse courriel :

\begin{Shaded}
\begin{Highlighting}[]
\FunctionTok{git}\NormalTok{ config }\AttributeTok{{-}{-}global}\NormalTok{ user.email }\StringTok{"votre.adresse@courriel.ca"}
\end{Highlighting}
\end{Shaded}

\begin{tcolorbox}[enhanced jigsaw, arc=.35mm, title=\textcolor{quarto-callout-note-color}{\faInfo}\hspace{0.5em}{Quelle adresse utiliser?}, colback=white, left=2mm, rightrule=.15mm, toptitle=1mm, opacityback=0, bottomtitle=1mm, breakable, colframe=quarto-callout-note-color-frame, coltitle=black, toprule=.15mm, titlerule=0mm, opacitybacktitle=0.6, bottomrule=.15mm, colbacktitle=quarto-callout-note-color!10!white, leftrule=.75mm]

Idéalement, utilisez une adresse courriel associée à votre compte
GitHub.

Cela permettra à GitHub d'associer correctement vos commits à votre
profil.

\end{tcolorbox}

Nous pouvons vérifier notre configuration :

\begin{Shaded}
\begin{Highlighting}[]
\FunctionTok{git}\NormalTok{ config }\AttributeTok{{-}{-}global}\NormalTok{ user.name}
\FunctionTok{git}\NormalTok{ config }\AttributeTok{{-}{-}global}\NormalTok{ user.email}
\end{Highlighting}
\end{Shaded}

Git devrait alors retourner les informations que nous venons d'entrer.

\bookmarksetup{startatroot}

\chapter{Vérification finale}\label{vuxe9rification-finale}

Avant de poursuivre, assurez-vous que les trois éléments suivants
fonctionnent.

Dans le terminal :

\begin{Shaded}
\begin{Highlighting}[]
\FunctionTok{git} \AttributeTok{{-}{-}version}
\end{Highlighting}
\end{Shaded}

doit retourner une version de Git.

Puis :

\begin{Shaded}
\begin{Highlighting}[]
\FunctionTok{git}\NormalTok{ config }\AttributeTok{{-}{-}global}\NormalTok{ user.name}
\end{Highlighting}
\end{Shaded}

et :

\begin{Shaded}
\begin{Highlighting}[]
\FunctionTok{git}\NormalTok{ config }\AttributeTok{{-}{-}global}\NormalTok{ user.email}
\end{Highlighting}
\end{Shaded}

doivent retourner vos informations.

Finalement, vous devez être capable de vous connecter à votre compte
\textbf{GitHub} dans votre navigateur.

\begin{tcolorbox}[enhanced jigsaw, arc=.35mm, title=\textcolor{quarto-callout-tip-color}{\faLightbulb}\hspace{0.5em}{Nous sommes prêts}, colback=white, left=2mm, rightrule=.15mm, toptitle=1mm, opacityback=0, bottomtitle=1mm, breakable, colframe=quarto-callout-tip-color-frame, coltitle=black, toprule=.15mm, titlerule=0mm, opacitybacktitle=0.6, bottomrule=.15mm, colbacktitle=quarto-callout-tip-color!10!white, leftrule=.75mm]

À ce stade :

\begin{Shaded}
\begin{Highlighting}[]
\NormalTok{Git est installé}
\NormalTok{       +}
\NormalTok{Git connaît notre identité}
\NormalTok{       +}
\NormalTok{notre compte GitHub existe}
\end{Highlighting}
\end{Shaded}

Nous pouvons maintenant connecter un \textbf{dépôt Git local} à un
\textbf{dépôt GitHub distant} et effectuer notre premier \texttt{push}.

\end{tcolorbox}

\section{Première partie du projet
(évaluation)}\label{premiuxe8re-partie-du-projet-uxe9valuation}

Une partie de la première évaluation consiste à crée votre compte GitHub
et y héberger votre librairie R!

\bookmarksetup{startatroot}

\chapter*{References}\label{references}
\addcontentsline{toc}{chapter}{References}

\markboth{References}{References}

\phantomsection\label{refs}
\begin{CSLReferences}{0}{1}
\end{CSLReferences}



\end{document}
